\documentclass{article}

\title{Ch. 31 HW \#3 Solutions}
\author{Principles of Macroeconomics}
\date{\today}

\begin{document}

\maketitle

\section*{Quick Quiz}

\begin{enumerate}

\setcounter{enumi}{5}

\item (c) -- If the price level is rising, that means the price of labor -- a.k.a. wage -- is rising as well.

\item (d) -- Since the inflation is consistent, people's inflation expectations will match the actual inflation. Arbitrary redistributions between debtors and creditors only occur when inflation expectations differ from actual inflation.

\item (c) -- Since loans are written in nominal terms, the real value of the payments depends on inflation. If inflation is high, the real value of the payments is lower, so creditors are hurt. 

\end{enumerate}

\section*{Questions for Review}

\begin{enumerate}

\setcounter{enumi}{5}

\item The costs of inflation include:

	\begin{itemize}

	\item Shoeleather costs
	
	\item Menu costs
	
	\item Distorted relative prices
	
	\item Increased taxes on saving
	
	\item A distorted unit of account
	
	\item Arbitrary shifts in wealth between creditors and debtors
	
	\end{itemize}
	
	Since inflation in the U.S. is usually low and predictable, most of these costs are minimized. The increased taxes on saving are probably the least mitigated and therefore the most important for the economy. Disincentivized saving makes our capital stock grow more slowly.

\item If inflation is less than expected, that means the real interest rate is larger than expected. That means the payments that debtors make to creditors are worth more, so creditors benefit.

\end{enumerate}

\section*{Problems \& Applications}

\begin{enumerate}

\setcounter{enumi}{4}

\item

	\begin{enumerate}
	
	\item Since the price of both goods doubled, the inflation rate was 100\%. Neither Bob nor Rita are better off. Although their costs doubled, so did their incomes because the price of the goods they sell doubled.
	
	\item Let's assume that the typical consumption basket is one of each good. Then the cost of the basket increased from \$4 to \$6. That means the inflation rate was 50\%.
	
	Bob is better off because his costs increased by 33\%, the percent change in the price of rise, but his income increased 100\%, the percent change in the price of beans.
	
	Rita is worse off because her costs increased by 100\%, but her income only increased by 33\%
	
	\item The cost of the basket decreased from \$4 to \$3.50. That means the inflation rate was
		\[ \frac{3.50 - 4}{4} = -12.5\%. \]
		
		Bob is better off because his costs decreased by 50\%, but his income increased by 100\%. 
		
		Rita is worse off because her costs increased by 100\%, but her income decreased by 50\%. 
		
	\item The relative price of rice and beans matters more to Bob and Rita than the overall inflation rate. We see in (a)-(c) that their well-being depends more on the relative prices than it does the overall inflation rate.
	
	\end{enumerate}

\end{enumerate}

\end{document}